\documentclass[11pt]{article}

% =====================================================================
%  The Detectability Spectrum of Imaginary-Quadratic Invariants
%  in Isomonodromy Data
%  STAGED DRAFT — NO deposit, NO submission. Conservative evidence classes.
% =====================================================================

\usepackage[utf8]{inputenc}
\usepackage{amsmath,amssymb,amsthm}
\usepackage{mathtools}
\usepackage[margin=1in]{geometry}
\usepackage{hyperref}
\usepackage{enumitem}

\newtheorem{theorem}{Theorem}
\newtheorem{lemma}[theorem]{Lemma}
\newtheorem{proposition}[theorem]{Proposition}
\newtheorem{corollary}[theorem]{Corollary}
\newtheorem{conjecture}[theorem]{Conjecture}
\theoremstyle{definition}
\newtheorem{definition}[theorem]{Definition}
\newtheorem{problem}[theorem]{Problem}
\newtheorem{remark}[theorem]{Remark}

\newcommand{\Q}{\mathbb{Q}}
\newcommand{\Z}{\mathbb{Z}}
\newcommand{\C}{\mathbb{C}}
\newcommand{\Cl}{\operatorname{Cl}}
\newcommand{\disc}{\operatorname{disc}}
\newcommand{\Res}{\operatorname{Res}}
\newcommand{\rank}{\operatorname{rank}}
\newcommand{\Kr}[2]{\left(\tfrac{#1}{#2}\right)}
\newcommand{\Dc}{\mathsf{D}}
\newcommand{\Tch}{\mathsf{T}}
\newcommand{\Mch}{\mathsf{M}}
\newcommand{\Stch}{\mathsf{S}}

% evidence-class tags
\newcommand{\PROVEN}{\textsf{[PROVEN]}}
\newcommand{\VERIFIED}{\textsf{[VERIFIED]}}
\newcommand{\STRUCTURAL}{\textsf{[STRUCTURAL]}}
\newcommand{\CONJ}{\textsf{[CONJECTURED]}}
\newcommand{\UNK}{\textsf{[UNKNOWN]}}
\newcommand{\VAC}{\textsf{[VACUOUS]}}

\title{The Detectability Spectrum of Imaginary-Quadratic Invariants\\
in Isomonodromy Data}
\author{(staged draft --- author/affiliation withheld pending operator)}
\date{Draft compiled \today}

\begin{document}
\maketitle

\begin{abstract}
For the self-adjoint degree-two polynomial-continued-fraction (PCF) family with
linear ODE $(a y')'-x^2y=0$, $a=\beta_2x^2+\beta_1x+\beta_0$, we organise the
question \emph{which arithmetic invariants of the imaginary quadratic field
$K=\Q(\sqrt{\Delta})$, $\Delta=\disc(a)$, leave a trace in the associated
Painlev\'e/isomonodromy data} into a \emph{detectability spectrum} indexed by
(invariant $\times$ channel) pairs, with channels surface type, monodromy, and
Stokes data. Our headline result is a \emph{$\Delta$-tower no-channel theorem}: under
the standard Sakai/Okamoto premise that the surface type is a function of the local
jet, the discrete Sakai surface type $T$ is selected by $\beta_0$-free local data and is
therefore constant on the locus $L=\{\beta_2,\beta_1,a_n\text{ fixed};\ \beta_0\
\text{free};\ \Delta\neq0\}$, so $T$ detects \emph{no} non-trivial arithmetic
invariant of $K$ --- not the class number $h$, the class group $\Cl(K)$, the unit
order $w$, the set of ramified primes, nor the splitting type of any fixed prime.
This strictly strengthens the previously recorded ``CM field, not class number''
observation from an empirical statement about one channel to a structural theorem
about the discrete invariant. We make precise the distinction between
\emph{recoverable-in-principle} (R) and \emph{intrinsic-channel-trace} (C)
detectability, express the spectrum through a functor $F$ from isomonodromy
structures to arithmetic invariants, and characterise $\ker F$: trivial-as-a-channel
for surface type (it forgets everything), and --- for the richer connection
channel --- trivial in sense (R) but conjecturally equal to the $L$-value stratum
$\{h,\Cl\}$ in sense (C). All symbolic and numerical claims are reproducible from a
single self-contained script. We separate cleanly what is proven, verified,
structural, and conjectural.
\end{abstract}

\tableofcontents

% =====================================================================
\section{Introduction}
% =====================================================================

Polynomial continued fractions (PCFs) attach to a recurrence
$P_n=b_nP_{n-1}+a_nP_{n-2}$ a second-order linear ODE whose isomonodromic
deformation carries a Sakai surface type, Stokes data at its irregular point, and a
monodromy representation. A recurring empirical theme in this corpus is that the
\emph{class number} of the imaginary quadratic field naturally attached to such a
family appears to be \emph{invisible}: in the degree-two case, six $\Delta<0$
families with class numbers spanning $h=1$ and $h=2$ exhibit identical Stokes
($S<1$) and PSLQ-non-detection behaviour, evidence that ``the controlling invariant
is the sign of $\Delta$ \ldots\ and not the class number'', sharpened by the pair
in the \emph{same} field $\Q(\sqrt{-7})$ at $h=1$ versus $h=2$ which identifies
``the CM field, not the class number'' as the structural carrier
\cite{PCF1}.

That observation is empirical and channel-specific. The present paper asks the
structural question in full generality. We fix the self-adjoint degree-two family
\begin{equation}\label{eq:ode}
  (a(x)\,y')'-x^2y=0,\qquad a(x)=\beta_2x^2+\beta_1x+\beta_0,\qquad
  b:=a',\quad c:=-x^2,
\end{equation}
with $\Delta=\disc(a)=\beta_1^2-4\beta_2\beta_0$ and $K=\Q(\sqrt{\Delta})$, and we
study \emph{every} standard arithmetic invariant of $K$ against \emph{each} of three
information channels. The organising object is a partition: each channel induces a
partition of the parameter space, and an arithmetic invariant is detectable through a
channel exactly when it is constant on the channel's fibres.

\paragraph{Contributions.}
\begin{enumerate}[leftmargin=2em]
\item A \emph{$\Delta$-tower no-channel theorem} (Theorem~\ref{thm:main}): the
discrete surface type $T$ is constant on $L$, hence detects no non-trivial arithmetic
invariant of $K$. The class-number case is the headline corollary; the proof is a
symbolic computation that every type-selector is $\beta_0$-free, and is reproducible
(\S\ref{sec:repro}, anchor module \texttt{nochannel/}).
\item A precise formalisation of the \emph{(R) vs.\ (C)} distinction
(\S\ref{sec:RC}): recoverable-in-principle versus intrinsic-channel-trace, the
conceptual spine that explains how the connection channel can recover the field $K$
yet leave the deep $L$-value invariants $\{h,\Cl\}$ with no known intrinsic reading.
\item A functor $F$ from isomonodromy structures to arithmetic invariants and a
characterisation of $\ker F$ per channel (\S\ref{sec:functor}).
\item A quick-results section (\S\ref{sec:quick}): the regulator is vacuous for
imaginary quadratic $K$; unit-order detectability reduces to a finite reachability
check of $\Delta\in\{-3,-4\}$; and the $\Delta$-tower invisibility is essentially
free from the main theorem.
\item Open problems (\S\ref{sec:open}): the intrinsic $h$-trace in monodromy/Stokes
data, with its conjectural link to $L(1,\chi_\Delta)$, and class-group structure
detection.
\end{enumerate}

\paragraph{Evidence discipline.} We tag every assertion with an evidence class:
\PROVEN\ (a complete argument, given clearly-stated standard premises),
\VERIFIED\ (exact machine computation), \STRUCTURAL\ (a precise statement supported
by the structure but without a written proof here), \CONJ, \UNK, and \VAC\ (vacuously
constant, no content). A reconciliation table (\S\ref{sec:reconcile}) maps each
numbered result to its \texttt{claims.jsonl} entry and flags any place where prose
risks over-stating the recorded class.

% =====================================================================
\section{Setup: channels, invariants, and the detectability map}\label{sec:setup}
% =====================================================================

\subsection{The parameter space and the field}
Fix $\beta_2>0$ and a numerator polynomial $a_n$ (the PCF data; $a_n\equiv1$ for the
V\_quad anchor, $a_n=n$ for the family QL15). The parameter space is
\[
  \mathcal P=\{p=(\beta_2,\beta_1,\beta_0)\in\Q^3:\ \Delta(p)\neq0\},\qquad
  \Delta(p)=\beta_1^2-4\beta_2\beta_0 .
\]
Each $p$ determines~\eqref{eq:ode}, its singular set $\{x_-,x_+,\infty\}$ with
$x_\pm=(-\beta_1\pm\sqrt{\Delta})/(2\beta_2)$, and the field $K(p)=\Q(\sqrt{\Delta(p)})$.

\begin{remark}[discriminant normalisation \STRUCTURAL]\label{rmk:disc}
We write $\Delta$ for $\disc(a)$. The field invariant is the squarefree core: $K$
depends only on the squarefree part of $\Delta$, and $\disc(a)$ may differ from the
\emph{fundamental} discriminant $d_K$ by a square. When $\Delta$ is itself a
fundamental discriminant (e.g.\ $\Delta=-20$ for QL15) the order $\mathcal O_\Delta$
is the maximal order and the form class number coincides with $h_K$. In general we
take $H(\Delta)$ to be the class number of the order $\mathcal O_\Delta$ (the number
of $SL_2(\Z)$-classes of primitive positive-definite forms of discriminant $\Delta$);
this is the invariant that varies along our loci.
\end{remark}

\subsection{Channels}
A \emph{channel} is a map $c\colon\mathcal P\to \Dc_c$ that factors through the
isomonodromic/gauge equivalence appropriate to its data type. We consider three.
\begin{itemize}[leftmargin=2em]
\item \textbf{Surface type} $\Tch\colon\mathcal P\to\{\text{Sakai surface types}\}$:
the discrete label of the space of initial conditions of the isomonodromic
deformation of~\eqref{eq:ode}, in the sense of Sakai \cite{Sakai2001} and Okamoto.
\item \textbf{Monodromy / connection} $\Mch$: the full Riemann--Hilbert datum --- the
linear monodromy representation, the connection (Stokes-completed) matrices, and the
\emph{positions} $x_\pm$ of the singular points --- modulo simultaneous conjugacy.
\item \textbf{Stokes} $\Stch$: the Stokes multipliers / constants at the rank-one
irregular point at $\infty$ (for the V\_quad anchor the relevant scalar is
$S=2\pi K$, the late-term resurgence constant).
\end{itemize}
These are listed in increasing fineness in the sense relevant below: $\Tch$ is a
coarse discrete invariant, $\Mch$ retains continuous moduli, and $\Stch$ refines the
local analytic data at $\infty$.

\subsection{Arithmetic invariants and detectability}
Let $\mathcal K$ be the set of imaginary quadratic fields. We consider the invariants
\[
  \Delta,\quad h=h_K,\quad \Cl(K),\quad w=\#\mathcal O_K^\times,\quad
  \mathrm{Ram}(K)=\{p:\ p\mid d_K\},\quad
  \mathrm{Spl}_p(K)=\Kr{d_K}{p},\quad R_K,
\]
each a map $\iota\colon\mathcal K\to S_\iota$. Composing with $K\colon\mathcal P\to
\mathcal K$ gives $\iota\circ K\colon\mathcal P\to S_\iota$.

\begin{definition}[detectability through a channel]\label{def:detect}
An arithmetic invariant $\iota$ is \emph{detectable through the channel $c$} if
$\iota\circ K$ factors through $c$: there exists $\varphi$ with
$\iota\circ K=\varphi\circ c$ on $\mathcal P$; equivalently $\iota\circ K$ is constant
on every fibre $c^{-1}(\delta)$. It is \emph{invisible through $c$} if no such
$\varphi$ exists, i.e.\ some fibre of $c$ carries at least two values of $\iota\circ K$.
\end{definition}

The \emph{detectability spectrum} is the table of (invariant, channel) pairs together
with the status of Definition~\ref{def:detect}. Section~\ref{sec:RC} refines
``factors through'' into the (R)/(C) dichotomy.

% =====================================================================
\section{The $\Delta$-tower no-channel theorem}\label{sec:main}
% =====================================================================

The engine is a symbolic fact about~\eqref{eq:ode}. All four parts are verified by
exact computation in the anchor module (\S\ref{sec:repro}).

\begin{lemma}[$\beta_0$-free local jet \PROVEN/\VERIFIED]\label{lem:jet}
For~\eqref{eq:ode} with $\Delta\neq0$:
\begin{enumerate}[label=(\roman*),leftmargin=2.2em]
\item $b=a'=\beta_1+2\beta_2x$ and $c=-x^2$ are polynomial identities; $b$ is
$\beta_0$-free.
\item At each finite root $x_0$ of $a$ the ODE $y''+(a'/a)y'+(c/a)y=0$ has a regular
singular point with indicial polynomial $r(r-1)+p_0r+q_0=r^2$, where
$p_0=\Res_{x_0}(a'/a)=1$ (because $a'/a=\frac{d}{dx}\log a$) and
$q_0=\lim_{x\to x_0}(x-x_0)^2(c/a)=0$. Hence the local exponents are $\{0,0\}$ at both
finite roots, independently of $\beta_0$; the resonance resolves to an
\emph{apparent} (no-logarithm) singularity, forced by $b=a'$ for all coefficients.
\item At $x=\infty$, $x\,(a'/a)\to2$ (so $a'/a\sim2/x$) and $c/a\to-1/\beta_2$; the
point is irregular of Poincar\'e rank $1$ with leading exponential rate
$\lambda=\pm1/\sqrt{\beta_2}$ ($=\pm\theta_\infty$, $\theta_\infty=2/\sqrt{\beta_2}$),
independent of $\beta_0$.
\item $a(0)=\beta_0$, so $x=0$ is an ordinary point iff $\beta_0\neq0$.
\end{enumerate}
Consequently the entire type-selecting jet --- pole orders, local/formal exponents,
and apparent-singularity conditions --- is independent of $\beta_0$; the only
$\beta_0$-dependence of the configuration is through the positions
$x_\pm=(-\beta_1\pm\sqrt\Delta)/(2\beta_2)$.
\end{lemma}

\begin{proof}[Proof (premise and computation)]
Parts (i)--(iv) are the exact symbolic computations of the anchor verifier
(\texttt{nochannel/surface\_type\_nochannel\_verify.py}, results recorded in
\texttt{surface\_type\_nochannel\_results.json}). The residue $p_0=1$ is the
logarithmic-derivative identity; $q_0=0$ because $c/a$ has only a simple pole at a
simple root; the indicial polynomial is then $r^2$. At $\infty$ the stated limits give
$\lambda^2=1/\beta_2$. The reduction ``surface type is a function of this local jet''
is the standard Sakai/Okamoto theory of spaces of initial conditions
\cite{Sakai2001}; we take it as the (clearly stated) premise.
\end{proof}

\begin{theorem}[$\Delta$-tower no-channel \PROVEN, given the Sakai premise]\label{thm:main}
Let $L=\{\beta_2,\beta_1,a_n\text{ fixed};\ \beta_0\in\Z;\ \Delta\neq0\}\subset
\mathcal P$. Then the surface type is constant on $L$,
\[
  \Tch(\beta_2,\beta_1,\beta_0,a_n)=\Tch(\beta_2,\beta_1,a_n)=:T_0 ,
\]
while $H(\Delta(\beta_0))$ takes at least two distinct values on $L$. Consequently
$\Tch$ does not factor through $H$, and more strongly: for \emph{every} arithmetic
invariant $\iota$ of $K$ that is non-constant on $\{K(p):p\in L\}$ --- in particular
$\iota\in\{\Delta,\,h,\,\Cl,\,w,\,\mathrm{Ram},\,\mathrm{Spl}_p\}$ --- $\iota$ is
\emph{invisible through $\Tch$}. The only arithmetic distinction $\Tch$ makes is the
non-arithmetic confluence dichotomy $\Delta=0$ vs.\ $\Delta\neq0$.
\end{theorem}

\begin{proof}
Constancy of $\Tch$ on $L$ is Lemma~\ref{lem:jet}. For the second clause it suffices
that $H$ is non-constant on $L$: by exact reduced-form count, on the worked locus
$\beta_2=3,\beta_1=1$ ($\Delta=1-12\beta_0$, all fundamental) one has
$H(\beta_0{=}1{:}10)=1,3,2,5,3,7,3,8,3,10$, already two distinct values at
$\beta_0=1,2$. If $\iota\circ K=\varphi\circ\Tch$ held on $L$ then $\iota\circ K$
would be constant on $L$ (as $\Tch$ is), contradicting non-constancy of any such
$\iota$; for $h$ this is immediate, and $\Delta,\Cl,w,\mathrm{Ram},\mathrm{Spl}_p$ are
non-constant along $L$ as well (\S\ref{sec:quick}). The confluence remark is
Lemma~\ref{lem:jet}(iv) together with $\Delta=0\Leftrightarrow x_+=x_-$.
\end{proof}

\begin{corollary}[class number; the headline \PROVEN]\label{cor:h}
$\Tch$ does not factor through the class number $h$. On the QL15 locus
$\beta_2=3,\beta_1=-2$ ($\Delta=4-12\beta_0$, which passes through the QL15 field at
$\beta_0=2$, $\Delta=-20=\disc\Q(\sqrt{-5})$, $h=2$) the type is constant while
$H$ takes the distinct values $\{1,2,3,4,6\}$ (these are order class numbers, some at
non-fundamental $\Delta$); the points $\beta_0=1$
($K=\Q(\sqrt{-2})$, $\Delta=-8$, $h=1$) and $\beta_0=2$ ($K=\Q(\sqrt{-5})$,
$\Delta=-20$, $h=2$) --- \emph{both fundamental}, so here $h$ is the field class
number $h_K$ --- lie on the same surface $T_0$ with $\Delta\neq0$ and different $h$.
The class-number conclusion thus rests only on fundamental discriminants; the
non-fundamental orders enter no part of the argument.
\end{corollary}

\begin{remark}[what is, and is not, used \PROVEN]
The proof uses only (a) constancy of $\Tch$ (Lemma~\ref{lem:jet}, symbolic) and (b)
a finite computation that $H$ takes $\ge2$ values. It does \emph{not} use the
identity of the constant surface $T_0$, nor any asymptotic growth of $H$; in
particular it carries no exposure to the ineffectivity of the Siegel lower bound.
The deposited determination $T_0=\Q$-isomorphism class $D_5^{(1)}$ (Painlev\'e~V,
symmetry $W(A_3^{(1)})$, $\delta=-\tfrac12$) is \VERIFIED\ but lies \emph{outside} the
proof's critical path.
\end{remark}

% =====================================================================
\section{The (R)/(C) distinction}\label{sec:RC}
% =====================================================================

Definition~\ref{def:detect} conflates two phenomena that Theorem~\ref{thm:main} forces
us to separate. The connection channel $\Mch$ retains the singular positions $x_\pm$,
hence $\sqrt\Delta$, hence the field $K$; so \emph{once $K$ is recovered, every
invariant of $K$ is computable}. That is a weak, almost tautological sense of
detection. The sharp content of a ``spectrum'' lives in a stronger sense.

\begin{definition}[recoverable-in-principle, (R)]\label{def:R}
$\iota$ is \emph{(R)-recoverable through $c$} if $\iota\circ K$ factors through $c$ in
the sense of Definition~\ref{def:detect}. Equivalently $c$ refines the partition of
$\mathcal P$ induced by $\iota\circ K$.
\end{definition}

\begin{definition}[intrinsic-channel-trace, (C)]\label{def:C}
Fix for each channel $c$ an \emph{admissible reading class} $\Phi_c$ of maps out of
$\Dc_c$: the features one is permitted to ``read off''. For $\Mch$ we take $\Phi_\Mch$
to be the maps given by $\overline\Q$-algebraic functions of the entries of the
monodromy/connection matrices and of the singular positions $x_\pm$ (the Riemann--Hilbert
coordinates), \emph{excluding} the arithmetic evaluation $K\mapsto\iota(K)$ itself. Then
$\iota$ is \emph{(C)-detectable through $c$} if $\iota\circ K=\varphi\circ c$ for some
$\varphi\in\Phi_c$.
\end{definition}

By construction $(\mathrm C)\Rightarrow(\mathrm R)$. The admissible class $\Phi_\Mch$
formalises ``read, do not solve'': a $(C)$-trace is a bounded-complexity feature of the
datum, not the result of recovering $K$ and then invoking an oracle for $h(K)$. The
boundary of $\Phi_\Mch$ is a modelling choice; the conjectures below are robust to any
reasonable choice because no reading of the monodromy datum is \emph{known} to produce
$h$ other than via $K$.

\begin{proposition}[field recovery \STRUCTURAL]\label{prop:Krec}
$K$ is (C)-detectable through $\Mch$: the unordered pair $\{x_-,x_+\}$ of singular
positions is part of the datum, and $\sqrt\Delta=\beta_2(x_+-x_-)$ is an algebraic
function of it, so $K=\Q(\sqrt\Delta)$ is read off. Consequently \emph{every}
invariant of $K$ is (R)-recoverable through $\Mch$.
\end{proposition}
This is \STRUCTURAL: the map ``RH datum $\mapsto\{x_\pm\}\mapsto\sqrt\Delta$'' is
forced by~\eqref{eq:ode}, but a fully written Riemann--Hilbert reconstruction theorem
for this family is not part of the present corpus.

\begin{remark}[the spectrum is a property of the pair, on one axis]
Combining Theorem~\ref{thm:main} and Proposition~\ref{prop:Krec}: for $\Tch$ even (R)
fails for every non-trivial $\iota$ (the strongest invisibility); for $\Mch$, (R)
holds for all $\iota$, but (C) splits the invariants into a \emph{geometric stratum}
$\{\Delta,\mathrm{Ram},\mathrm{Spl}_p,w\}$ that are algebraic functions of $x_\pm$
(hence (C)-detectable) and an \emph{$L$-value stratum} $\{h,\Cl\}$ for which no
admissible reading is known (\S\ref{sec:open}).
\end{remark}

% =====================================================================
\section{The functor $F$ and its kernel}\label{sec:functor}
% =====================================================================

Let $\mathcal S$ be the groupoid of isomonodromy structures of the
family~\eqref{eq:ode} (objects: ODEs on $\mathcal P$; isomorphisms: isomonodromic
deformations and gauge equivalences). A channel $c$ is a functor $\mathcal S\to
\Dc_c$ to a discrete (for $\Tch$) or moduli (for $\Mch$) target; an invariant
$\iota$ is a functor $\mathcal S\to S_\iota$ via $K$. Write $F_c$ for the induced map
on objects and define the \emph{indistinguishability kernel}
$\ker F_c=\{(p,p'):c(p)=c(p')\}$, the equivalence relation collapsed by $c$.

\begin{proposition}[surface-type functor \PROVEN, given the Sakai premise]\label{prop:kerT}
$F_{\Tch}$ is constant on each locus $L$ (Theorem~\ref{thm:main}, under the standard
Sakai/Okamoto premise that $\Tch$ is a function of the local jet), so
$\ker F_{\Tch}\supseteq L\times L$.
Hence the composite $\iota\circ K$ is constant on $L$ for the partition induced by
$\Tch$, and the entire arithmetic spectrum
$\{\Delta,h,\Cl,w,\mathrm{Ram},\mathrm{Spl}_p\}$ lies in the part of $\mathcal P$ that
$F_{\Tch}$ cannot separate. In words: \emph{the discrete isomonodromy invariant is
arithmetically blind} --- it detects no invariant of $K$ beyond the non-arithmetic
dichotomy $\Delta=0/\!\neq\!0$.
\end{proposition}

\begin{proposition}[connection functor \STRUCTURAL/\CONJ]\label{prop:kerM}
$F_{\Mch}$ factors through $K$ (Proposition~\ref{prop:Krec}). Therefore:
\begin{itemize}[leftmargin=2em]
\item the \emph{provable} kernel is trivial in sense (R): $\ker F_{\Mch}$ separates any
two parameters with $K(p)\neq K(p')$, so no invariant of $K$ lies in $\ker F_{\Mch}$
(every invariant is a function of the recovered $K$); \STRUCTURAL;
\item the \emph{effective} kernel in sense (C) --- the invariants with no admissible
reading $\varphi\in\Phi_\Mch$ beyond ``recover $K$, then evaluate'' --- contains the
$L$-value stratum: $\ker^{(C)}F_{\Mch}\supseteq\{h,\Cl\}$ is the content of
Conjecture~\ref{conj:h}. \CONJ
\end{itemize}
\end{proposition}

Thus $\ker F$ has a clean two-tier description: for the discrete channel it swallows
the whole arithmetic spectrum; for the connection channel it is trivial in (R) and
conjecturally exactly the $L$-value stratum in (C). This is the strongest general
statement the corpus currently supports.

% =====================================================================
\section{Quick results}\label{sec:quick}
% =====================================================================

\subsection{The regulator is vacuous \VAC/\PROVEN}
For every imaginary quadratic $K$ the unit group has rank $0$, so by the Dirichlet
unit theorem the regulator is the empty determinant $R_K=1$. A constant function
carries no information: ``invisibility'' of $R_K$ is vacuous, not a no-channel
phenomenon. Content appears only for \emph{real} quadratic fields, where
$R_K=\log\varepsilon_0$ varies; we note that QL15's field $\Q(\sqrt{-5})$ is
imaginary, so this axis is dormant here and marks the natural next domain.

\subsection{Ramification, splitting, and the geometric stratum \PROVEN(to $\Tch$)/\STRUCTURAL(to $\Mch$)}
$\mathrm{Ram}(K)=\{p\mid d_K\}$ and $\mathrm{Spl}_p(K)=\Kr{d_K}{p}$ are functions of
$\Delta$. By Theorem~\ref{thm:main} they are invisible through $\Tch$ \PROVEN. Through
$\Mch$ they are (C)-detectable as algebraic functions of $\sqrt\Delta$
(Proposition~\ref{prop:Krec}) \STRUCTURAL. We record the arithmetic nuance: the full
splitting pattern over \emph{all} primes equals the character $\chi_\Delta$, which
determines $\Delta$ (a primitive real Dirichlet character determines its modulus), so
all-prime splitting is equivalent to $K$; but finitely many small primes
under-determine $K$.

\subsection{Unit order and the reachability of $\Delta\in\{-3,-4\}$ \VERIFIED}
The unit order satisfies $w=6\iff\Delta=-3$, $w=4\iff\Delta=-4$, else $w=2$; thus $w$
is a function of $\Delta$ and is invisible through $\Tch$ \PROVEN. Whether the
extra-symmetry fields $\Q(i),\Q(\sqrt{-3})$ (with CM $j=1728,0$) are realised on a
given locus is a finite check: on $a=x^2+\beta_0$ one reaches $\Delta=-4$ at
$\beta_0=1$; on $a=3x^2+3x+\beta_0$ one reaches $\Delta=-3$ at $\beta_0=1$; but on the
QL15 locus $a=3x^2-2x+\beta_0$ one reaches \emph{neither} for integer $\beta_0$, since
$\Delta=-4$ requires $3\beta_0-1$ to be a perfect square, impossible as $3\beta_0-1
\equiv2\pmod 3$. \VERIFIED\ (reproduce, \S\ref{sec:repro}). Whether $w>2$ leaves an
extra-symmetry trace in $\Mch$ is \STRUCTURAL.

% =====================================================================
\section{Open problems}\label{sec:open}
% =====================================================================

\begin{conjecture}[no intrinsic $h$-trace \CONJ]\label{conj:h}
$h$ is (C)-invisible through $\Mch$ and through $\Stch$: there is no admissible reading
$\varphi\in\Phi_\Mch$ (resp.\ $\Phi_\Stch$) with $h=\varphi\circ\Mch$ on the family.
Equivalently, $h$ is recoverable only via the field $K$, not as a bounded-complexity
feature of the monodromy or Stokes datum.
\end{conjecture}
The heuristic is the class-number formula
$h=\dfrac{w\sqrt{|d_K|}}{2\pi}\,L(1,\chi_\Delta)$: $h$ is a \emph{global} $L$-value,
whereas $\Phi_\Mch$ comprises local algebraic features. A first concrete probe is to
sweep the Stokes constant $S=2\pi K$ across the $\beta_0$-family and test for any
correlation with $h$; currently only a single data point ($S$ at the V\_quad anchor)
is computed, so the Stokes case is \UNK\ rather than \CONJ. Proving
Conjecture~\ref{conj:h} (invisibility) would require a theorem that \emph{no} algebraic
function of the RH coordinates equals $h$ along the family; exhibiting a trace
(detectability) would require constructing the channel --- e.g.\ relating an invariant
trace field or a monodromy-group arithmetic invariant to $L(1,\chi_\Delta)$. Either
direction needs machinery beyond the present corpus.

\begin{problem}[class-group structure \UNK]\label{prob:Cl}
Detect $\Cl(K)$ (the group, not merely $h=\#\Cl$) through $\Mch$ or $\Stch$. This is
strictly harder than Conjecture~\ref{conj:h}: it asks for the group structure, and is
recoverable in sense (R) via $K$ but has no known intrinsic reading.
\end{problem}

\begin{problem}[real-quadratic extension \STRUCTURAL]
Repeat the analysis for parameter regimes with $\Delta>0$, where the regulator
acquires content ($R_K=\log\varepsilon_0$) and the unit group becomes infinite. The
self-adjoint structure $b=a'$ persists, so Lemma~\ref{lem:jet} carries over and the
surface-type invisibility should hold verbatim; the new question is whether the
regulator or fundamental unit leaves a trace in $\Mch/\Stch$.
\end{problem}

% =====================================================================
\section{The spectrum, summarised}\label{sec:summary}
% =====================================================================

\begin{center}\small
\begin{tabular}{l l l l}
\hline
Invariant & vs.\ $\Tch$ & vs.\ $\Mch$ & vs.\ $\Stch$\\
\hline
$\Delta$ (field $K$)     & invisible \PROVEN              & detectable \STRUCTURAL & moduli-dep.\ \STRUCTURAL\\
$\mathrm{Ram}$, $\mathrm{Spl}_p$ & invisible \PROVEN     & detectable \STRUCTURAL & \UNK\\
unit order $w$           & invisible \PROVEN             & extra-symmetry \STRUCTURAL & \UNK\\
class number $h$         & invisible \PROVEN             & (R) yes / (C) \CONJ\ no    & \UNK\\
class group $\Cl$        & invisible \PROVEN\ (a fortiori)& (R) yes / (C) \UNK         & \UNK\\
regulator $R_K$          & vacuous \VAC                  & vacuous \VAC               & vacuous \VAC\\
\hline
\end{tabular}
\end{center}

\noindent In the ``vs.\ $\Tch$'' column, every \PROVEN\ entry is \PROVEN\ \emph{modulo}
the standard Sakai/Okamoto premise of Theorem~\ref{thm:main} (surface type a function
of the local jet); none is asserted unconditionally.

\noindent The one-axis reading: a \emph{geometric stratum} ($\Delta$ and its
elementary functions) is invisible to the discrete type but faithful to $K$ in the
connection moduli; an \emph{$L$-value stratum} ($h,\Cl$) is invisible to the discrete
type and conjecturally has no intrinsic trace anywhere; the regulator is vacuous for
imaginary quadratic $K$.

% =====================================================================
\section{Reconciliation with the claims ledger}\label{sec:reconcile}
% =====================================================================

Every numbered result maps to a recorded claim; we flag any prose that might be read
above its recorded class.

\begin{center}\small
\begin{tabular}{l l l l}
\hline
Result & Ledger entry & Recorded class & Note\\
\hline
Lemma~\ref{lem:jet}        & \texttt{nochannel} NC-2..NC-5 & \VERIFIED\ (symbolic) & exact sympy\\
Theorem~\ref{thm:main}     & \texttt{nochannel} NC-7       & \PROVEN$^\ast$        & $\ast$ given Sakai premise\\
Corollary~\ref{cor:h}      & \texttt{nochannel} NC-6, DS-H-ST & \PROVEN/\VERIFIED  & QL15 locus, $h(-20)=2$\\
Prop.~\ref{prop:Krec}      & DS-DISC-MON                   & \STRUCTURAL           & not a written RH theorem\\
Prop.~\ref{prop:kerT}      & DS-FUNCTOR-T                  & \PROVEN               & from Thm~\ref{thm:main}\\
Prop.~\ref{prop:kerM}      & DS-FUNCTOR-M                  & \STRUCTURAL/\CONJ     & (R) trivial, (C) conj.\\
\S\ref{sec:quick} regulator& DS-REG-ALL                    & \VAC/\PROVEN          & $R_K\equiv1$\\
\S\ref{sec:quick} $w$ reach.& DS-UNIT-MON                  & \VERIFIED             & finite check\\
Conj.~\ref{conj:h}         & DS-H-MON, DS-H-STK            & \CONJ/\UNK            & $L$-value heuristic only\\
Prob.~\ref{prob:Cl}        & DS-CL-MON                     & \UNK                  & harder than Conj.~\ref{conj:h}\\
\hline
\end{tabular}
\end{center}

\paragraph{Flags (prose vs.\ ledger).}
\begin{itemize}[leftmargin=2em]
\item Theorem~\ref{thm:main} is \PROVEN\ \emph{only modulo} the Sakai/Okamoto premise
that $\Tch$ is a function of the local jet; the qualifier is stated at every
occurrence of the result --- abstract, contribution list (via back-reference to
Theorem~\ref{thm:main}), the theorem and its proof, Proposition~\ref{prop:kerT}, and
the summary table note. Do not cite Theorem~\ref{thm:main} as unconditional.
\item Proposition~\ref{prop:Krec} (``$\Mch$ recovers $K$'') is \STRUCTURAL, not
\PROVEN: it relies on a Riemann--Hilbert reconstruction that is asserted, not written.
Any downstream use of ``$\Mch$ detects $\Delta$'' inherits \STRUCTURAL.
\item The phrase ``$\Mch$ detects everything'' (sense (R)) must never be softened into
``$\Mch$ detects $h$ intrinsically''; the latter is exactly Conjecture~\ref{conj:h},
recorded \CONJ.
\item The surface label $D_5^{(1)}$ is \VERIFIED\ (deposited, MEDIUM band) and is kept
\emph{off} the critical path; the paper must not present it as first-principles
\PROVEN.
\end{itemize}

% =====================================================================
\section{Reproducibility}\label{sec:repro}
% =====================================================================

All symbolic and numerical claims are reproduced by a single self-contained script
(CAS: \texttt{sympy}, exact symbolic for the ODE algebra; exact integer
form-counts for class numbers, cross-checked against twelve classical anchors
including $h(-20)=2$, $h(-3)=h(-4)=1$ with unit orders $6,4$).

\begin{verbatim}
# foundational anchor (surface-type no-channel, Lemma 3.1 / Theorem 3.2):
python nochannel/surface_type_nochannel_verify.py
# detectability spectrum numerics (class numbers, splitting, reachability):
python files/detectability/detect_spectrum_verify.py
\end{verbatim}

The anchor verifier exits $0$ iff: the indicial polynomial is $r^2$ at both finite
roots, $\lambda=\pm1/\sqrt{\beta_2}$ at $\infty$ (all $\beta_0$-free); the QL15-locus
class numbers are $\{1,2,3,4,6\}$ with $\Delta=-20,h=2$ at $\beta_0=2$; and the
prior-anchor sequence $1,3,2,5,3,7,3,8,3,10$ is reproduced exactly.

% =====================================================================
\section*{Venue scan (editorial note, not for publication)}
% =====================================================================
Three candidate venues consistent with the corpus history, \emph{excluding}
Experimental Mathematics (blacklisted) and \emph{Notes on Number Theory and Discrete
Mathematics} (blocked):
\begin{enumerate}[leftmargin=2em]
\item \textbf{SIGMA (Symmetry, Integrability and Geometry: Methods and Applications)}
--- open access; strong on Painlev\'e/isomonodromy and Sakai-surface classification;
the surface-type framing fits squarely. Foreground the isomonodromy structure.
\item \textbf{The Ramanujan Journal} --- continued fractions and arithmetic; the
PCF/class-number framing and computer-assisted evidence are within scope. Foreground
the no-channel theorem as a continued-fraction/arithmetic statement.
\item \textbf{Journal of Physics A: Mathematical and Theoretical} --- isomonodromy,
Stokes phenomena, Painlev\'e; the channel/Stokes framing and open problems fit.
Foreground the detectability spectrum and the $L$-value open problem.
\end{enumerate}
Fit caveats: SIGMA and J.\ Phys.\ A expect an integrable-systems narrative; the
Ramanujan Journal expects the arithmetic angle foregrounded. The result is structural
and conservative, which suits all three.

% =====================================================================
\begin{thebibliography}{9}
\bibitem{PCF1} \emph{A CM predicate for degree-two polynomial continued fractions}
(PCF-1), corpus manuscript \texttt{p12\_pcf1\_main.tex}; the empirical
``sign of $\Delta$ / CM field, not the class number'' statement and the
$\Q(\sqrt{-7})$ pair (QL06, QL26).
\bibitem{Sakai2001} H.~Sakai, \emph{Rational surfaces associated with affine root
systems and geometry of the Painlev\'e equations}, Comm.\ Math.\ Phys.\ \textbf{220}
(2001), 165--229.
\bibitem{Takeuchi1977} K.~Takeuchi, \emph{Arithmetic triangle groups}, J.\ Math.\
Soc.\ Japan \textbf{29} (1977), 91--106.
\bibitem{Vquad} V\_quad anchor, Zenodo \texttt{10.5281/zenodo.20455090}
(Painlev\'e~V on Sakai $D_5^{(1)}$; verdict VQ-N1).
\end{thebibliography}

\end{document}
